\documentclass[11pt]{article}
\usepackage[margin=1in]{geometry}
\usepackage{amsmath,amssymb,amsthm,booktabs,array,longtable,microtype}
\usepackage[colorlinks=true,linkcolor=blue,citecolor=blue,urlcolor=blue]{hyperref}
\newtheorem{theorem}{Theorem}
\newtheorem{lemma}[theorem]{Lemma}
\newtheorem{corollary}[theorem]{Corollary}
\DeclareMathOperator{\lcm}{lcm}
\title{Proving Hanna's conjecture in OEIS A181074, Comments (2 October 2010): extensions of Moree (2003) and Gossow (2024)}
\author{Haobo Ma\\{\small ChronoAI Pte Ltd, Singapore}\\[0.7em]
Wenlin Zhang\\{\small National University of Singapore}\\{\small ChronoAI Pte Ltd, Singapore}}
\date{8 November 2026\\[0.5em]{\small Creative Commons Attribution 4.0 International (CC BY 4.0)}}
\begin{document}
\maketitle
\begin{abstract}
Extending coefficient congruences of Moree (2003) and Gossow (2024),
we prove Paul D. Hanna's conjecture in the Comments of OEIS A181074
(entry dated 2 October 2010) that the coefficients of
\(\exp(\sum_{n\ge1,k\ge0}\binom{n+k-1}{k}^{k+1}x^{n+k}/n)\)
are integers \cite{moree,gossow,h74}. The proof lifts a coefficient congruence modulo a prime to
prime powers and applies the Euler-product criterion for integral formal
exponentials. The same argument proves the integrality conjectures in
A181070, A181076, A181078, A181080, A181082 and A181084, and gives an integral-series
theorem for the coefficients of powers of any series in \(1+t\mathbb Z[[t]]\).
The reflected exponent also extends to every integer polynomial with constant
coefficient one. The companion logarithmic coefficients satisfy Dold congruences.
A bivariate Euler product proves integrality of the triangles A209424,
A228899, A209196 and A228904 and their row sums.
An entropy estimate proves Kotesovec's limit formulas in A209428, Formula
(6 March 2014), and A166895, Formula (20 November 2024).
Three exponent-lowered variants have nonintegral
coefficients already in degree six, as asserted in the source Comments.
\end{abstract}

\section{The conjecture and a general statement}

Hanna's entry A181074 defines
\begin{equation}\label{eq:main}
 A(x)=\exp\left(\sum_{n\ge1}\sum_{k\ge0}
     \binom{n+k-1}{k}^{k+1}\frac{x^{n+k}}{n}\right)
\end{equation}
and its Comments state: ``Conjecture: this sequence consists entirely of
integers'' \cite{h74}. We prove this assertion for every coefficient, together
with its directly related binomial variants. All series below are formal:
in degree \(m\), the sums involve only \(1\le n\le m\) and \(k=m-n\).
No analytic convergence is required.

Moree's integral Witt transform \cite[Theorem 4(2)]{moree}, first released
in 2003, already gives the unpowered coefficient-array divisibility.
Gossow \cite[Theorem 3.3, (3.1), and Example 3.4]{gossow} places
$[t^k]U(t)^n$ and both binomial arrays in the general theory of Gauss
congruences on ranked semigroups. We extend these coefficient congruences
to the index-dependent powers below; closure under a fixed product does
not alone give those varying exponents. Our coefficient-lifting lemma
and bivariate Euler-product criterion are reproved antecedents, not new
general criteria. Du, Huang and Li \cite[Theorem 2(2)--(3)]{dhl} give the
one-variable exponential/Euler-product equivalence in 2005; Byszewski,
Graff and Ward \cite[Theorem 2.5 and Lemma 2.6]{bgw} give its later Dold-sequence
formulation. We give the elementary arguments needed for these extensions.

\begin{theorem}\label{thm:general}
Let \(U(t)\in1+t\mathbb Z[[t]]\), and write
\(c(n,k)=[t^k]U(t)^n\) for \(n\ge1\), \(k\ge0\).
Let the nonnegative exponent $E(n,k)$ have one of the forms
\begin{align*}
 E(n,k)&=rk+s &&(r,s\ge1),\\
 E(n,k)&=rn+sk+q &&(r\ge1,\ s,q\ge0),\\
 E(n,k)&=n+k-1.
\end{align*}
All parameters here are integers. Then
the series
\begin{equation}\label{eq:general}
 A_{U,E}(x)=\exp\left(\sum_{n\ge1,k\ge0}
                      c(n,k)^{E(n,k)}\frac{x^{n+k}}{n}\right)
\end{equation}
belongs to \(1+x\mathbb Z[[x]]\).
\end{theorem}

The exponent zero occurs only for \(E=n+k-1\), \(n=1\), \(k=0\);
its base is \(c(1,0)=1\). The theorem thus involves no ambiguous power.
Taking \(U=(1-t)^{-1}\) and \(E=k+1\) gives \eqref{eq:main}.

\section{Two coefficient congruences}

The following unpowered lemma also follows from Moree's Witt-transform
integrality and Gossow's coefficient-array example \cite{moree,gossow}.

For a prime \(p\) and nonzero integer \(z\), let \(v_p(z)\) be its
\(p\)-adic valuation, with the same convention for rational numbers.

\begin{lemma}\label{lem:coeff}
For \(U\in\mathbb Z[[t]]\), \(n\ge1\), \(k\ge0\), and prime \(p\),
\begin{align}
 n&\mid k\,c(n,k),\label{eq:div}\\
 c(pn,pk)&\equiv c(n,k)\pmod {p^{v_p(n)+1}}.\label{eq:lift}
\end{align}
\end{lemma}
\begin{proof}
The identity \(t(U^n)'=ntU^{n-1}U'\) proves \eqref{eq:div} by taking
the coefficient of \(t^k\), including \(k=0\).
Reduction modulo \(p\) gives \(U(t)^p\equiv U(t^p)\pmod p\).
This identity applies to formal series coefficient by coefficient, since
each coefficient depends on a finite truncation of \(U\).

If \(V\equiv W\pmod {p^r}\) in \(\mathbb Z[[t]]\), \(r\ge1\), the
binomial expansion gives \(V^p\equiv W^p\pmod {p^{r+1}}\): the intermediate
binomial coefficients are divisible by \(p\), and \(rp\ge r+1\) controls
the last term. This includes \(p=2\).
Write \(n=p^a b\), where \(a=v_p(n)\). Raise the preceding congruence
between \(U^p\) and \(U(t^p)\) to the power \(p^a\), then to the power
\(b\). It follows that
\[
 U(t)^{pn}\equiv U(t^p)^n\pmod {p^{a+1}}.
\]
Taking the coefficient of \(t^{pk}\) proves \eqref{eq:lift}.
\end{proof}

\begin{lemma}\label{lem:powered}
Under the hypotheses of Theorem~\ref{thm:general}, put
\(F(n,k)=c(n,k)^{E(n,k)}\). Then
\begin{equation}\label{eq:array}
 n\mid kF(n,k),\qquad
 F(pn,pk)\equiv F(n,k)\pmod {p^{v_p(n)+1}}
\end{equation}
for every \(n\ge1\), \(k\ge0\), and prime \(p\).
\end{lemma}
\begin{proof}
The first assertion follows from \eqref{eq:div} whenever \(E\ge1\).
When \(E=0\), we have \(k=0\), so the assertion still holds.
For the second assertion let \(a=v_p(n)\), \(c=c(n,k)\),
\(e=E(n,k)\) and \(e'=E(pn,pk)\).
By \eqref{eq:lift}, it suffices to show
\(c^{e'}\equiv c^e\pmod {p^{a+1}}\).
If \(k=0\), then \(c=c(pn,0)=1\), so this is immediate.
For \(k>0\), both exponents are positive, and \(e'\ge e\).
The case \(c=0\) is therefore immediate as well.

Suppose first that \(p\nmid c\). From \(n\mid kc\) we obtain
\(p^a\mid k\). For the three exponent families, the difference \(e'-e\) is respectively
\[
 (p-1)rk,\qquad (p-1)(rn+sk),\qquad (p-1)(n+k).
\]
In each case it is divisible by \(p^a(p-1)=\varphi(p^{a+1})\).
Euler's theorem modulo \(p^{a+1}\) proves the desired congruence.

Now suppose that \(p\mid c\), and set \(b=v_p(c)\ge1\).
Equation~\eqref{eq:div} implies \(a\le v_p(k)+b\).
For \(e=rk+s\), we have $e\ge k+1$. The inequality \(v_p(k)\le k-1\) gives
\[
 eb\ge(k+1)b\ge v_p(k)+b+1\ge a+1.
\]
For the other two families, \(k>0\) gives \(e\ge n\), while
\(n\ge p^a\ge a+1\), so again \(eb\ge a+1\).
Both powers are thus zero modulo \(p^{a+1}\).
\end{proof}

\section{From an array to an integral exponential}

\begin{lemma}\label{lem:criterion}
Let \(F(n,k)\in\mathbb Z\), \(n\ge1\), \(k\ge0\), satisfy
\eqref{eq:array}. Define
\begin{equation}\label{eq:ghost}
 g_m=\sum_{n=1}^m\frac{m}{n}F(n,m-n),\qquad m\ge1.
\end{equation}
Then \(g_m\in\mathbb Z\), and for every prime \(p\), \(m\ge1\),
\begin{equation}\label{eq:dold}
 g_{pm}\equiv g_m\pmod {p^{v_p(m)+1}}.
\end{equation}
Consequently
\begin{equation}\label{eq:euler}
 \exp\left(\sum_{m\ge1}g_m\frac{x^m}{m}\right)
  =\prod_{m\ge1}(1-x^m)^{-h_m}\in1+x\mathbb Z[[x]],
 \qquad h_m=\frac1m\sum_{d\mid m}\mu(d)g_{m/d}\in\mathbb Z.
\end{equation}
\end{lemma}
\begin{proof}
Each summand in \eqref{eq:ghost} is integral, since, with \(k=m-n\),
\[
 \frac{m}{n}F(n,k)=F(n,k)+\frac{kF(n,k)}{n}.
\]
To compare \(g_{pm}\) with \(g_m\), split its summation index \(n\)
according to divisibility by \(p\), and put \(b=v_p(m)\).
If \(p\nmid n\), its summand \(pmF(n,pm-n)/n\) is an integer of
valuation at least \(b+1\), hence vanishes modulo \(p^{b+1}\).
For \(n=pj\), pair the summand with index \(j\) in \(g_m\).
Their difference is the integer
\[
 \frac{m}{j}\bigl(F(pj,p(m-j))-F(j,m-j)\bigr).
\]
By \eqref{eq:array} its valuation is at least
\(v_p(m)-v_p(j)+v_p(j)+1=b+1\). This proves \eqref{eq:dold}.

For completeness, let \(p^a\Vert m\), and write \(m=p^a r\), \(p\nmid r\).
The numerator defining \(h_m\) equals
\[
 \sum_{d\mid r}\mu(d)
       \bigl(g_{m/d}-g_{m/(pd)}\bigr).
\]
Every difference is divisible by \(p^a\), by \eqref{eq:dold} applied to
\(m/(pd)\). This holds for every prime power dividing \(m\), so \(h_m\)
is an integer. M\"obius inversion gives \(g_m=\sum_{d\mid m}dh_d\).
Taking formal logarithms now proves the identity in \eqref{eq:euler}.
Each factor has integral coefficients, even when \(h_m<0\), and only
finitely many factors affect a fixed degree. The product is therefore integral.
\end{proof}

\begin{proof}[Proof of Theorem~\ref{thm:general}]
Lemma~\ref{lem:powered} supplies \eqref{eq:array}.
Grouping \eqref{eq:general} by \(m=n+k\) identifies its logarithm with
\(\sum_{m\ge1}g_m x^m/m\). Lemma~\ref{lem:criterion} finishes the proof.
\end{proof}

\section{Binomial specializations and source coverage}

Take \(U=(1-t)^{-1}\) or \(U=1+t\). Then respectively
\[
 c(n,k)=\binom{n+k-1}{k},\qquad c(n,k)=\binom nk,
\]
where the second coefficient is zero for \(k>n\).
The following reflected-polynomial extension supplies A181080 as well.

\begin{corollary}\label{cor:polynomial}
Let $P(t)\in1+t\mathbb Z[t]$ have degree at most $d\ge1$, and let
$r,s\ge1$ be integers. Put $c(n,k)=[t^k]P(t)^n$. Then
\[
 \exp\left(\sum_{n\ge1}\sum_{k=0}^{dn}
 c(n,k)^{r(dn-k)+s}\frac{x^{n+k}}n\right)\in1+x\mathbb Z[[x]].
\]
\end{corollary}
\begin{proof}
Extend the powered array by zero for $k>dn$. Its divisibility follows
from $n\mid kc(n,k)$. To establish its coefficient congruence, put
$a=v_p(n)$, $c=c(n,k)$ and $j=dn-k$. The support is preserved by $(n,k)\mapsto(pn,pk)$,
and $n\mid jc(n,k)$ because $n\mid kc(n,k)$.
The coefficient lift in Lemma~\ref{lem:coeff} reduces the assertion to
$c^{rpj+s}\equiv c^{rj+s}\pmod{p^{a+1}}$.
If $j=0$, the exponents are equal. If $j>0$ and $p\nmid c$,
then $p^a\mid j$, so Euler's theorem applies to the exponent difference
$(p-1)rj$. If $p\mid c\ne0$, write $b=v_p(c)$.
Then $a\le v_p(j)+b$ and
$(rj+s)b\ge(j+1)b\ge v_p(j)+b+1\ge a+1$.
Both powers vanish modulo $p^{a+1}$; $c=0$ is immediate.
Lemma~\ref{lem:criterion} proves the assertion.
\end{proof}

\begin{corollary}\label{cor:binomial}
The eight exponential series in Table~\ref{tab:coverage} have nonnegative
integer coefficients.
\end{corollary}
\begin{proof}
For A181074 and A181070 take $E=k+1$ in Theorem~\ref{thm:general}.
For A181076 and A166894 take $E=n$; for A181078 take $E=n+k-1$;
for A181082 and A181084 take $E=n+k$ and $E=n+k+1$.
For A181080 use Corollary~\ref{cor:polynomial} with $P=1+t$ and
$d=r=s=1$. Each logarithm has nonnegative coefficients, so its
exponential does too.
\end{proof}

\begin{table}[ht]
\centering\small
\begin{tabular}{@{}llll@{}}
\toprule
OEIS entry & Base & Exponent & Coverage of displayed statement\\
\midrule
A181074 & $\binom{n+k-1}{k}$ & $k+1$ & Conjecture proved\\
A181076 & $\binom{n+k-1}{k}$ & $n$ & Conjecture proved\\
A181078 & $\binom{n+k-1}{k}$ & $n+k-1$ & Conjecture proved\\
A181070 & $\binom nk$ & $k+1$ & Conjecture proved\\
A181080 & $\binom nk$ & $n-k+1$ & Conjecture proved\\
A181082 & $\binom nk$ & $n+k$ & Conjecture proved\\
A181084 & $\binom nk$ & $n+k+1$ & Conjecture proved\\
A166894 & $\binom nk$ & $n$ & Integrality stated in Name; proved\\
\bottomrule
\end{tabular}
\caption{The exponentials are defined by summing $F(n,k)x^{n+k}/n$.
For ordinary binomial bases the range is $0\le k\le n$.
The first seven entries label integrality a conjecture in their Comments;
A166894 states it in its Name.}\label{tab:coverage}
\end{table}

These locators are the Name and Comments fields of Hanna's entries
\cite{h74,h70,h76,h78,h80,h82,h84,h166}. In all these binomial cases the
logarithm has nonnegative coefficients, so the integral exponential does too.
The directly linked A181075 is also covered: \eqref{eq:ghost} gives
\begin{equation}\label{eq:log}
 g_m=\sum_{k=0}^{m-1}\binom{m-1}{k}^{k+1}\frac{m}{m-k},
 \qquad xA'(x)/A(x)=\sum_{m\ge1}g_mx^m.
\end{equation}
Thus its defining sum is integral and satisfies \eqref{eq:dold}.
The other logarithmic identities follow by the same substitution $n=m-k$.
For completeness their coverage is listed below \cite{logs}.
\begin{corollary}\label{cor:logs}
Each of the eight logarithmic sequences below is integral and satisfies
the Dold congruences \eqref{eq:dold}.
\end{corollary}
\begin{proof}
For each exponential, its powered array satisfies \eqref{eq:array}, as
proved in Theorem~\ref{thm:general} or Corollary~\ref{cor:polynomial}.
Apply Lemma~\ref{lem:criterion} and substitute $n=m-k$ in
\eqref{eq:ghost}; logarithmic differentiation identifies the displayed sum.
\end{proof}
\begin{center}\small
\begin{tabular}{@{}lll@{}}
\toprule
Logarithmic entry & Exponential entry & Exponent after $n=m-k$\\
\midrule
A181075 & A181074 & $k+1$\\
A181077 & A181076 & $m-k$\\
A181079 & A181078 & $m-1$\\
A181071 & A181070 & $k+1$\\
A181081 & A181080 & $m-2k+1$\\
A181083 & A181082 & $m$\\
A181085 & A181084 & $m+1$\\
A166895 & A166894 & $m-k$\\
\bottomrule
\end{tabular}
\end{center}
The range is $0\le k\le m-1$ for the first three rows and
$0\le k\le\lfloor m/2\rfloor$ for the remaining rows.
The base is respectively $\binom{m-1}{k}$ or $\binom{m-k}{k}$,
and each summand has the additional factor $m/(m-k)$.
The separate limit in A166895 is established in the next section.

Three cross-reference details can be resolved directly from the formulas.
A181075's cross-reference names A181076 as its exponential, whereas its
displayed sum is \eqref{eq:log} for A181074. A181076 instead has logarithmic
coefficients A181077. A181079's Formula field also names A181076,
but its defining sum is the logarithmic derivative of A181078.
The displayed Name of A181081 has a stray multiplication
sign in its exponent; differentiation of A181080 gives the unambiguous exponent
\(m-2k+1\) above. The A181072 link in A181070 leads to a random-walk entry;
its displayed statement is not a binomial-exponential problem and is not
part of this coverage \cite{walk}. Likewise, the A181073 links in A181071
and A181077 lead to a different entry about digit permutations, not a
binomial-exponential sequence \cite{digits}.

\section{The companion limit and coverage boundaries}

V. Kotesovec gives the following limit for the unweighted sum in A209428,
Formula (6 March 2014), and for the weighted sum in A166895, Formula
(20 November 2024) \cite{limit209,logs}. Both entries state the limit
without a proof. Its root is identified separately in A323777 \cite{root}.
\begin{corollary}\label{cor:limit}
For $g_m=\sum_{k=0}^{\lfloor m/2\rfloor}
\binom{m-k}{k}^{m-k}m/(m-k)$, one has
\[
 \lim_{m\to\infty}g_m^{1/m^2}
 =\left(\frac{1-r}{r}\right)^{(1-r)^2/(3-4r)},
\]
where $r$ is the unique root in $(0,1/2)$ of
$(1-2r)^{3-4r}=(1-r)^{2(1-r)}r^{1-2r}$.
The same limit holds for the defining sum
$b_m=\sum_{k=0}^{\lfloor m/2\rfloor}\binom{m-k}{k}^{m-k}$ of A209428.
\end{corollary}
\begin{proof}
Use $0\log0=0$ and define, for $0\le u\le1/2$,
\[
 H(u)=(1-u)\log(1-u)-u\log u-(1-2u)\log(1-2u),
 \qquad f(u)=(1-u)H(u).
\]
The elementary binomial bounds
\[
 \frac{e^{Nh(j/N)}}{N+1}\le\binom Nj\le e^{Nh(j/N)},
 \qquad h(v)=-v\log v-(1-v)\log(1-v),
\]
hold for $0\le j\le N$. Indeed, in the binomial probability
distribution with parameter $j/N$ the $j$th term is a maximum,
between $1/(N+1)$ and $1$; endpoints follow directly.
Apply these bounds with $N=m-k$ and $j=k$.
The logarithm of each summand of $g_m$ is
$m^2f(k/m)+O(m\log(m+1))$, uniformly in $k$.
The factor $m/(m-k)$ lies between $1$ and $2$, and there are at most
$m+1$ terms. Consequently
$m^{-2}\log g_m\to\max_{0\le u\le1/2}f(u)$,
by continuity and the density of these grids.

In the interior,
\[
 H'(u)=2\log(1-2u)-\log u-\log(1-u),\qquad
 H''(u)=-\frac{4}{1-2u}-\frac1u+\frac1{1-u}<0.
\]
Thus $H'$ decreases from $+\infty$ to $-\infty$.
Where $H'>0$, $f''=-2H'+(1-u)H''<0$; where $H'\le0$,
$f'=-H+(1-u)H'<0$, since $H>0$ in the interior.
As $f'(0+)=+\infty$, $f$ has exactly one maximizer $r$.
Expanding $f'(r)=0$ gives the stated equation, and substitution gives
\[
 f(r)=\frac{(1-r)^2}{3-4r}\log\frac{1-r}{r}.
\]
Exponentiating proves the formula for $g_m$. Finally,
$b_m\le g_m\le2b_m$, so taking $m^{-2}$th powers proves the same limit
for $b_m$.
\end{proof}

The coverage above concerns the displayed mathematical assertions, not
an integrality conjecture for the lowered exponents: the source explicitly
states their nonintegrality. The following additional coverage rows make
this distinction precise.
\begin{center}\small
\begin{tabular}{@{}>{\raggedright\arraybackslash}p{.40\linewidth}>{\raggedright\arraybackslash}p{.54\linewidth}@{}}
\toprule
Source item and locator & Coverage\\
\midrule
A181071, A181075, A181077, A181079, A181081, A181083, A181085, A166895, Name
 & Proved: integral defining sums and Dold congruences, Corollary~\ref{cor:logs}.\\
A166895, Formula, Kotesovec (20 November 2024)
 & Proved: the stated limit, Corollary~\ref{cor:limit}; credited to Kotesovec.\\
A209428, Formula, Kotesovec (6 March 2014)
 & Proved: the companion limit, Corollary~\ref{cor:limit}.\\
A323777, Name, defining root
 & Not posed; uniqueness in $(0,1/2)$ is proved in Corollary~\ref{cor:limit}.\\
General reflected-polynomial series, Corollary~\ref{cor:polynomial}
 & Not posed; integrality is proved for $\deg P\le d$, $r,s\ge1$.\\
A181074, A181070, A181080, Comments, lowered powers
 & Proved: the asserted nonintegrality. Integrality is not asserted in the
 source; it fails in degree six as below.\\
A181075, Crossrefs; A181079, Formula; A181081, Name
 & Display/cross-reference corrections explained above; no conjecture is refuted.\\
A181072, Name and Formula; A181073, Name and Comments
 & Not covered: the random-walk and digit-permutation entries concern
 different sequences.\\
\bottomrule
\end{tabular}
\end{center}

\section{Bivariate triangles and their specializations}

Hanna's A209424, Name (8 March 2012), and A228899, Name
(7 September 2013), define triangles by retaining a second variable in
the ordinary-binomial exponentials \cite{triangles}. The general
Gauss-congruence criterion of Gossow \cite[Theorem 3.3]{gossow}
has the following Euler-product form.

\begin{corollary}\label{cor:bivariate}
Suppose $F(n,k)\in\mathbb Z$ satisfies \eqref{eq:array} for $n\ge1$, $k\ge0$,
and for each $n$ has finite support in $k$. Then
\[
 B_F(x,y)=\exp\left(\sum_{n\ge1,k\ge0}F(n,k)\frac{x^ny^k}{n}\right)
 =\prod_{n\ge1,k\ge0}(1-x^ny^k)^{-h_{n,k}}
 \in1+x\mathbb Z[y][[x]],
\]
where $h_{n,k}=n^{-1}\sum_{d\mid\gcd(n,k)}\mu(d)F(n/d,k/d)$.
In particular the four triangles in Table~\ref{tab:bivariate} are integral.
\end{corollary}
\begin{proof}
Fix $p^a\Vert n$. If $p\nmid k$, every divisor $d$ of $\gcd(n,k)$
is prime to $p$, and $(n/d)\mid(k/d)F(n/d,k/d)$ makes each term
divisible by $p^a$. If $p\mid k$, pair squarefree divisors $d,pd$,
with $p\nmid d$. Their difference is divisible by $p^a$ by the second
congruence in \eqref{eq:array}. This also applies to $k=0$, with
$\gcd(n,0)=n$. Hence $h_{n,k}\in\mathbb Z$. M\"obius inversion gives
$F(n,k)=\sum_{d\mid\gcd(n,k)}(n/d)h_{n/d,k/d}$, so taking formal
logarithms proves the product. Each $x$-degree involves finitely many
$y$-degrees by the support assumption, proving the claimed ring membership.

For the first two rows of Table~\ref{tab:bivariate},
Lemma~\ref{lem:powered} supplies \eqref{eq:array}. For the last two,
write $(N,K)=(n^2,nk)$ or $(nk,k^2)$. The binomial identity
$N\mid K\binom NK$ implies $n\mid kF(n,k)$; for the second pair
cancel $k>0$, with $k=0$ immediate. Under $(n,k)\mapsto(pn,pk)$,
both $N,K$ are multiplied by $p^2$. Two applications of
\eqref{eq:lift} to $\binom NK=[t^K](1+t)^N$ give the required
modulus $p^{v_p(n)+1}$, since $v_p(N)\ge v_p(n)$.
For the second pair $k=0$ has $F(n,0)=1$ directly.
\end{proof}

\begin{table}[ht]\centering\small
\begin{tabular}{@{}lll@{}}
\toprule
Triangle (Name) & $F(n,k)$, $0\le k\le n$ & Coverage\\
\midrule
A209424 & $\binom nk^n$ & Proved: integral triangle\\
A228899 & $\binom nk^{k+1}$ & Proved: integral triangle\\
A209196 & $\binom{n^2}{nk}$ & Proved: integral triangle\\
A228904 & $\binom{nk}{k^2}$ & Proved: integral triangle\\
\bottomrule
\end{tabular}
\caption{Integrality of the displayed triangles, Corollary~\ref{cor:bivariate}.
Each source presents an integer triangle, not an explicit integrality
conjecture. The last two entries are dated 5 March 2012 and
7 September 2013.}\label{tab:bivariate}
\end{table}

\begin{corollary}\label{cor:rows}
The row sums A167007 of A209424 and A184730 of A228899, and the
column A209425 and central terms A209426 of A209424, are integral.
Their antidiagonal sums are respectively A166894 and A181070.
Column one of A209424 and the subdiagonal of A228899 both equal
$\sum_{j=1}^n j^{j-1}$, the defining formula of A060946.
\end{corollary}
\begin{proof}
Apply $y=1$ for row sums and $(x,y)=(z,z)$ for antidiagonal sums to
Corollary~\ref{cor:bivariate}. A column or central subsequence consists
of individual integer coefficients. These are the defining identities
in the source entries \cite{triangles,rows}.
For column one, extract the coefficient of $y$ from
$(1-x)^{-1}\exp(\sum_{n\ge1,k\ge1}\binom nk^nx^ny^k/n)$,
which is $(1-x)^{-1}\sum_{n\ge1}n^{n-1}x^n$.
For the subdiagonal replace $B(x,y)$ by $B(xy,y^{-1})$;
the exponent becomes $n-k+1$, giving the same coefficient of $y$.
Compare A060946, Formula \cite{trace}.
\end{proof}

\begin{corollary}\label{cor:lowered-triangles}
The lowered-power variants asserted nonintegral in A228899 and A184730,
Comments, are nonintegral.
\end{corollary}
\begin{proof}
For A228899's lowered-power variant, its Comments already assert
nonintegrality. Its $y=1$ specialization is the variant asserted
nonintegral in A184730, Comments. Exact expansion gives
$[x^4y^2]B_{\binom nk^k}=31/2$ and
$[x^4]B_{\binom nk^k}(x,1)=85/2$, proving both assertions.
Neither is a refuted integrality conjecture.
The coefficients follow from the exact recurrence
$n[x^ny^k]B=\sum_{j=1}^n\sum_{i=0}^j\binom ji^i
[x^{n-j}y^{k-i}]B$, with $[x^0]B=1$.
\end{proof}

\begin{corollary}\label{cor:arrays}
The arrays A209427 and A219207 are integral and satisfy
\eqref{eq:array}. A maximizing column $k_n$ of row $n$ of A219207 satisfies
$k_n/n\to u$, where $u$ is the unique root in $(1/2,1)$ of
$(1-u)^{2u-1}=u^{2u}$.
\end{corollary}
\begin{proof}
The logarithmic arrays A209427 and A219207 are, by their Names,
$\binom nk^n$ and $\binom nk^{k+1}$; their integrality follows
directly. Their Gauss congruences follow from Lemma~\ref{lem:powered}.
A219207's additional maximal-position assertion is also elementary:
the binomial bounds in Corollary~\ref{cor:limit} give the uniform limit
$n^{-2}\log\binom n{\lfloor un\rfloor}^{\lfloor un\rfloor+1}
\to u h(u)$ on $[0,1]$. Its derivative is positive on $(0,1/2]$;
on $(1/2,1)$ its second derivative
$2\log((1-u)/u)-1/(1-u)$ is negative, and its first derivative tends
to $-\infty$. Thus the maximizer is unique and satisfies
$(1-u)^{2u-1}=u^{2u}$, so maximizing indices divided by $n$
converge to that root, as asserted by Kotesovec in A219207, Comments
(15 November 2012) \cite{arrays}.
\end{proof}

\begin{center}\small
\begin{tabular}{@{}>{\raggedright\arraybackslash}p{.40\linewidth}>{\raggedright\arraybackslash}p{.54\linewidth}@{}}
\toprule
Source item and locator & Coverage\\
\midrule
Moree (2003), Theorem 4(2); Gossow (2024), Example 3.4
 & Reproved; prior proofs Moree (2003), Gossow (2024): unpowered congruences.\\
Du--Huang--Li (2005), Theorem 2; Gossow (2024), Theorem 3.3
 & Reproved; prior proofs cited: exponential/Euler-product criteria.\\
A209424, A228899, A209196, A228904, Name
 & Proved: stated integer triangles, Corollary~\ref{cor:bivariate}.\\
A167007, A184730, Name; A209425, A209426, Name
 & Proved: integral specializations, Corollary~\ref{cor:rows}.\\
A228899 and A184730, Comments, lowered powers
 & Proved: stated nonintegrality by exact rational coefficients.\\
A209427, A219207, Name; A219207, Comments (15 November 2012)
 & Proved: integral arrays, their congruences, and the stated maximizing-position limit.\\
A060946, Formula, $\sum_{j=1}^n j^{j-1}$
 & Prior formula cited \cite{trace}; its integrality is immediate.\\
\bottomrule
\end{tabular}
\end{center}

\section{The lowered exponents and exact checks}

The Comments of A181074, A181070 and A181080 also state that replacing
\(k+1\) by \(k\), or \(n-k+1\) by \(n-k\), does not give an integer
series. This is confirmed by exact counterexamples in degree six:
\begin{center}\small
\begin{tabular}{@{}lll@{}}
\toprule
Base and exponent & Coefficients in degrees $0,\ldots,5$ & $[x^6]$\\
\midrule
$\binom{n+k-1}{k}^{k}$ & $1,1,2,4,11,59$ & $1481/2$\\
$\binom nk^{k}$ & $1,1,2,3,5,10$ & $45/2$\\
$\binom nk^{n-k}$ & $1,1,2,3,7,25$ & $327/2$\\
\bottomrule
\end{tabular}
\end{center}
Here the finite range for the ordinary binomial sums remains \(0\le k\le n\).
For a direct rational calculation, their respective ghost coefficients
\((g_1,\ldots,g_6)\) are
\[
 (1,3,7,27,231,4038),\quad
 (1,3,4,7,21,63),\quad
 (1,3,4,15,86,807).
\]
Starting with \(a_0=1\), the identity
\begin{equation}\label{eq:recurrence}
 ma_m=\sum_{i=1}^m g_i a_{m-i}
\end{equation}
then gives each displayed fraction. This provides explicit counterexamples
to integrality of the lowered variants, in agreement with the source Comments.

For A181074, \((g_1,\ldots,g_6)=(1,3,10,71,1026,30912)\), and the first
Euler exponents in \eqref{eq:euler} are \(1,1,3,17,205,5150\).
Its coefficients begin
\[
 1,1,2,5,23,231,5405,322799,42761356,12597156231,9136063939651.
\]
The accompanying Python program uses exact rational arithmetic in
\eqref{eq:recurrence} and independently expands \eqref{eq:euler}.
It matches all 76 values in Greubel's A181074 table \cite{table74}, checks
integrality and both expansions through degree 75 for A181074 and degree 40
for the seven other exponentials in Table~\ref{tab:coverage}, and compares the
available displayed logarithmic coefficients. It also checks 5,544 coefficient
lifts and 29,664 powered congruences for \(1\le n\le24\), \(0\le k\le24\),
with primes \(2,3,5,7,11,13\). For ordinary binomials these local checks use
\(k\le n\). These are finite checks; the all-index conclusions follow from
the written proofs above. The reflected-polynomial extension is also checked
by both expansions through degree 24 for three integer polynomials,
including negative coefficients and a degree bound larger than the actual
degree. These finite checks do not establish the asymptotic limit, whose
proof is Corollary~\ref{cor:limit}.
For the four bivariate triangles, an exact polynomial recurrence and a
separate monomial Euler-product expansion agree through row 12, including
the displayed OEIS prefixes, row sums, column and central subsequences.
The two additional lowered-power fractions are checked in the same program.

\paragraph{Lean declarations and mathematical scope.}
At revision \texttt{abb48f7} (the full revision is fixed in each link), the declaration
\texttt{result} in
\href{https://github.com/the-omega-institute/trureturing/blob/abb48f723fe221da235c625800ba0d2ca1755296/D5/S1/Recurrence/Algebraic/HannaA181074Integrality.lean}{the A181074 Lean source}
has type \texttt{claim}, defined as
\(\forall m\in\mathbb N,\ \exists z\in\mathbb Z,\ [x^m]A=z\).
The declarations \texttt{result} in
\href{https://github.com/the-omega-institute/trureturing/blob/abb48f723fe221da235c625800ba0d2ca1755296/D5/S1/Recurrence/Algebraic/HannaA181070Integrality.lean}{the A181070 Lean source}
and
\href{https://github.com/the-omega-institute/trureturing/blob/abb48f723fe221da235c625800ba0d2ca1755296/D5/S1/Recurrence/Algebraic/HannaA181080Integrality.lean}{the A181080 Lean source}
have the corresponding coefficient-integrality statements for A181070 and
A181080. The general theorem for \(U\), the reflected-polynomial extension,
the other specializations, the logarithmic identities and the two companion limits
are supplied here as written proofs. None of those additional statements
is the conclusion of these three \texttt{result} declarations.
The bivariate Euler product, its triangle and subsequence corollaries,
and the A219207 maximizing-position limit likewise have written proofs.
The program's coverage is the finite range just stated.

\begin{thebibliography}{99}
\bibitem{moree} P. Moree, \emph{The formal series Witt transform},
arXiv:math/0311194, first released 12 November 2003; v2, 31 March 2005,
Theorem 4(2). \url{https://arxiv.org/abs/math/0311194v2}.
\bibitem{gossow} F. Gossow,
\emph{Polynomial and combinatorial analogues of Gauss congruence},
arXiv:2410.05678v2, 22 December 2024,
Theorem 3.3, (3.1), and Example 3.4.
\url{https://arxiv.org/abs/2410.05678v2}.
\bibitem{dhl} B.-S. Du, S.-S. Huang and M.-C. Li,
\emph{Newton, Fermat, and Exactly Realizable Sequences},
Journal of Integer Sequences \textbf{8} (2005), Article 05.1.2,
Theorem 2(2)--(3).
\url{https://cs.uwaterloo.ca/journals/JIS/VOL8/Huang/huang11.pdf}.
\bibitem{h74} P. D. Hanna, OEIS A181074, Name and Comments, 2 October 2010,
entry revision 9 (3 August 2024). \url{https://oeis.org/A181074}.
\bibitem{bgw} J. Byszewski, G. Graff and T. Ward,
\emph{Dold sequences, periodic points, and dynamics},
Bulletin of the London Mathematical Society \textbf{53} (2021), 1263--1298.
\href{https://doi.org/10.1112/blms.12531}{doi:10.1112/blms.12531}.
Accessible text: \href{https://arxiv.org/abs/2007.04031v4}{arXiv:2007.04031v4},
Theorem 2.5 and Lemma 2.6.
\bibitem{h70} P. D. Hanna, OEIS A181070, Name and Comments, 2 October 2010.
\url{https://oeis.org/A181070}.
\bibitem{h76} P. D. Hanna, OEIS A181076, Name and Comments, 2 October 2010.
\url{https://oeis.org/A181076}.
\bibitem{h78} P. D. Hanna, OEIS A181078, Name and Comments, 3 October 2010.
\url{https://oeis.org/A181078}.
\bibitem{h80} P. D. Hanna, OEIS A181080, Name and Comments, 2 October 2010.
\url{https://oeis.org/A181080}.
\bibitem{h82} P. D. Hanna, OEIS A181082, Name and Comments, 2 October 2010.
\url{https://oeis.org/A181082}.
\bibitem{h84} P. D. Hanna, OEIS A181084, Name and Comments, 28 October 2010.
\url{https://oeis.org/A181084}.
\bibitem{h166} P. D. Hanna, OEIS A166894, Name, 23 November 2009.
\url{https://oeis.org/A166894}.
\bibitem{logs} P. D. Hanna, OEIS A181071, A181075, A181077, A181081 and A181083
(2 October 2010), A181079 (3 October 2010), A181085 (28 October 2010)
and A166895 (23 November 2009), Name fields.
\url{https://oeis.org/A181071}; \url{https://oeis.org/A181075};
\url{https://oeis.org/A181077}; \url{https://oeis.org/A181079};
\url{https://oeis.org/A181081}; \url{https://oeis.org/A181083};
\url{https://oeis.org/A181085}; \url{https://oeis.org/A166895}.
\bibitem{walk} S. Perepechko, OEIS A181072, Name, 23 January 2011.
\url{https://oeis.org/A181072}.
\bibitem{digits} M. Rip\`a, OEIS A181073, Name and Comments, 23 January 2011.
\url{https://oeis.org/A181073}.
\bibitem{limit209} P. D. Hanna, OEIS A209428, Name (8 March 2012);
V. Kotesovec, limit in Formula (6 March 2014).
\url{https://oeis.org/A209428}.
\bibitem{root} V. Kotesovec, OEIS A323777, Name, 27 January 2019.
\url{https://oeis.org/A323777}.
\bibitem{table74} G. C. Greubel, Table of \(n,a(n)\), \(n=0,\ldots,75\),
linked from OEIS A181074. \url{https://oeis.org/A181074/b181074.txt}.
\bibitem{triangles} P. D. Hanna, OEIS A209424, Name and Comments
(8 March 2012); A228899, Name and Comments (7 September 2013);
A209196, Name (5 March 2012); A228904, Name (7 September 2013).
\url{https://oeis.org/A209424}; \url{https://oeis.org/A228899};
\url{https://oeis.org/A209196}; \url{https://oeis.org/A228904}.
\bibitem{rows} P. D. Hanna, OEIS A167007, Name (17 November 2009);
A184730, Name and Comments (20 January 2011);
A209425 and A209426, Name (8 March 2012).
\url{https://oeis.org/A167007}; \url{https://oeis.org/A184730};
\url{https://oeis.org/A209425}; \url{https://oeis.org/A209426}.
\bibitem{arrays} P. D. Hanna, OEIS A209427, Name (8 March 2012),
and A219207, Name (14 November 2012); V. Kotesovec, A219207,
Comments (15 November 2012). \url{https://oeis.org/A209427};
\url{https://oeis.org/A219207}.
\bibitem{trace} A. Fares, OEIS A060946, Name and Formula (8 May 2001).
\url{https://oeis.org/A060946}.
\end{thebibliography}
\end{document}
